\documentclass[11pt]{article}
\usepackage{amsmath,amssymb,amsthm,hyperref}
\newtheorem{theorem}{Theorem}
\newtheorem{lemma}{Lemma}
\begin{document}
\title{Large exact fibers of positive reciprocal heads on the Boolean cube}
\author{}\date{}\maketitle

\paragraph{Status.}
The results below are complete elementary written proofs, accompanied by
exact finite checks. They have not been formalized in Lean. They concern
the number of distinct representations of one real function, not the
number of sign functions that a head budget represents. In particular
they establish neither accepted asymptotic target. No model definition
or admissibility condition is changed.

\begin{theorem}[Histogram fibers]\label{thm:fibers}
Let $k\geq1$, $H\geq1$, and $n\geq3k$ be integers. There is a real
function $F$ on the Boolean $n$-cube having at least
\[
 \frac{(H!)^k}{[e(2k+1)]^{H7^k}}
\]
distinct unordered representations as a sum of $H$ reciprocals of
affine forms with constant coefficient one and strictly positive slopes.
Every such reciprocal is an original-model atom with constant numerator
one. A lower bound smaller than one is interpreted literally and is
uninformative.
\end{theorem}

\begin{proof}
Index the heads by $h\in\{0,\ldots,H-1\}$. For each $1\leq\ell\leq k$
choose a permutation $\pi_\ell$ of this set, independently. Put
\[
 a_\ell(x)=x_{3\ell-2}-x_{3\ell-1},\qquad
 b_\ell(x)=x_{3\ell-1}-x_{3\ell},\qquad
 D(x)=1+(H+1)\sum_{i=1}^{3k}x_i
       +(H+1)\sum_{j=1}^{n-3k}2^{j-1}x_{3k+j},
\]
and define
\[
 B_h^\pi(x)=D(x)+
       \sum_{\ell=1}^k\bigl(h a_\ell(x)+\pi_\ell(h)b_\ell(x)\bigr).
\]
The three slopes contributed within triple $\ell$, in addition to
the common slope $H+1$, are $h$, $\pi_\ell(h)-h$, and
$-\pi_\ell(h)$. Each final slope is at least two. Every unused
coordinate has a strictly positive slope. The constant is one, so all denominators
are positive on the whole cube. The existing conversion
\texttt{FracAtom.ofPositiveAffineRatio} realizes their reciprocals
as original-model atoms.

For a Boolean triple the pair $(a_\ell,b_\ell)$ belongs to
\[
 \mathcal V=\{(0,0),(1,0),(-1,0),(0,1),(0,-1),(1,-1),(-1,1)\}.
\]
For $v=((a_1,b_1),\ldots,(a_k,b_k))\in\mathcal V^k$, associate to
$\pi$ the multiset
\[
 M_v(\pi)=
 \left\{\sum_{\ell=1}^k[h a_\ell+\pi_\ell(h)b_\ell]:
                                  0\leq h<H\right\}.
\]
Its entries are integers in $[-k(H-1),k(H-1)]$. There are at most
\[
 \binom{H+2k(H-1)}{H}
 \leq\binom{(2k+1)H}{H}
 \leq[e(2k+1)]^H
\]
possible multisets. The last inequality follows from
$\binom{M}{H}\leq M^H/H!$ and $H!\geq(H/e)^H$.
Consequently the joint signature $(M_v)_{v\in\mathcal V^k}$ has at
most $[e(2k+1)]^{H7^k}$ possibilities. Among the $(H!)^k$ permutation
tuples, one signature has at least the asserted number of preimages.

For any two tuples in that fiber and any Boolean point $x$, the
corresponding multisets $\{B_h^\pi(x):h<H\}$ agree: they are
$M_{v(x)}(\pi)$ translated by the common number $D(x)$.
Therefore the reciprocal sums agree exactly at every Boolean point.

Finally distinct tuples give distinct unordered multisets of affine
forms. Indeed the coefficient of $x_1$ in $B_h^\pi$ is $H+1+h$,
so it uniquely identifies $h$. The coefficient of $x_{3\ell}$ then
recovers $\pi_\ell(h)$. Thus head relabeling does not divide the
fiber count by an additional factorial.
\end{proof}

\begin{lemma}[One triple]\label{lem:one}
For $k=1$, the joint signature is determined by the difference
multiset $\{h-\pi(h):h<H\}$. Hence the bound in
Theorem~\ref{thm:fibers} improves to $H!/2^{3H}$.
\end{lemma}
\begin{proof}
The seven multisets are zero, the fixed positive and negative index
sets, the same index sets permuted, and the difference multiset
and its negative. There are
$\binom{3H-2}{H}\leq2^{3H}$ possible difference multisets, since
each consists of $H$ entries in $\{-(H-1),\ldots,H-1\}$.
The same pigeonhole and distinctness arguments apply.
\end{proof}

\begin{theorem}[Superfactorial fibers in growing dimension]
Fix $\alpha>0$. Let $H=H(n)$ be any integer sequence with
$\log H\geq\alpha n$ for all sufficiently large $n$. There are
constants $c_\alpha>0$ and $N_\alpha$ such that for $n\geq N_\alpha$
one real function on the $n$-cube has at least
\[
 \exp(c_\alpha n H\log n)
\]
distinct unordered positive reciprocal representations of length $H$.
\end{theorem}
\begin{proof}
Take $k=\lfloor(\log n)/(2\log7)\rfloor$. For all sufficiently
large $n$, we have $3k\leq n$, $k\geq(\log n)/(3\log7)$, and
$7^k\leq\sqrt n$. The logarithm of the bound in
Theorem~\ref{thm:fibers} is at least
\[
 kH(\log H-1)-H7^k\log(e(2k+1)).
\]
The first term is at least
$\alpha nH\log n/(6\log7)$ eventually. The second is
$O(H\sqrt n\log\log n)=o(nH\log n)$. Increasing the threshold
therefore gives, for example,
$c_\alpha=\alpha/(12\log7)$.
\end{proof}

\begin{lemma}[Cauchy powers]\label{lem:cauchy}
Let $p\geq1$ be an integer and let $\lambda_1,\ldots,\lambda_H$
be distinct real numbers. On any interval where every
$t+\lambda_h>0$, a nonzero linear combination
$\sum_h c_h/(t+\lambda_h)^p$ has at most $H-1$ distinct zeros.
In particular its evaluations at any $H$ distinct points form an
invertible square matrix.
\end{lemma}
\begin{proof}
Induct on $H$, simultaneously for all positive integer exponents.
The assertion for one term is immediate. Multiply a putative
combination with $H$ distinct zeros by $(t+\lambda_1)^p$ and
differentiate. The derivative is
\[
 p(t+\lambda_1)^{p-1}
       \sum_{h=2}^H
       \frac{c_h(\lambda_h-\lambda_1)}{(t+\lambda_h)^{p+1}}.
\]
If some $c_h$ for $h>1$ is nonzero, the induction hypothesis bounds
the number of derivative zeros by $H-2$, contradicting Rolle's
theorem. If all those coefficients vanish, the original function
has no zero. The same induction also proves that no nonzero
coefficient vector represents the zero function. Singularity of
the evaluation matrix would supply exactly the forbidden linear
combination.
\end{proof}

\begin{theorem}[Regularity in a specified affine family]\label{thm:regular}
Suppose additionally that $2^{n-3k}\geq H$. All representations in
the fiber of Theorem~\ref{thm:fibers} lie in one common
$(k+1)H$-parameter affine family of denominator coefficients.
There is a fixed selection of $(k+1)H$ cube outputs such that the
Jacobian from those parameters to the selected reciprocal-sum
outputs is invertible at every permutation point. In particular
every displayed representation is locally isolated in the value
fiber within this specified parameter family.
\end{theorem}
\begin{proof}
The parameters are $d_h$ and $p_{h\ell}$, and the denominators are
\[
 B_h(x;d,p)=D(x)+d_h+
     \sum_{\ell=1}^k[h a_\ell(x)+p_{h\ell}b_\ell(x)].
\]
At a permutation point, $d_h=0$ and $p_{h\ell}=\pi_\ell(h)$.
Every such point has an open parameter neighborhood retaining
positive denominator constants, positive slopes, and positivity
on the cube.

Choose $H$ distinct filler vectors $y$; their contributions to $D$
are distinct by the binary filler weights. For each chosen $y$
use $k+1$ core patterns. The first has
$(a_1,b_1)=(1,0)$ and all other pairs zero. For each
$1\leq\ell\leq k$, the additional pattern has
$(a_\ell,b_\ell)=(0,1)$ and all other pairs zero. These are
actual Boolean patterns: use $100$ for $(1,0)$, $110$ for
$(0,1)$, and $000$ for zero. The first group of $H$ outputs
has derivative zero in every $p$ coordinate, and its derivative
in the $d$ coordinates is
\[
 -\frac{1}{(T_y+h)^2},
\]
where $T_y=D(100,000,\ldots,000,y)$ takes $H$ distinct positive
values. This block is invertible by Lemma~\ref{lem:cauchy}.

For the additional pattern indexed by $\ell$, the only nonzero
$p$ derivative block is the one for $p_{h\ell}$, and it equals
\[
 -\frac{1}{(T'_y+\pi_\ell(h))^2}.
\]
The $d$ block has the same entries, but lies strictly below the
first diagonal block. Every displayed diagonal block is
invertible by Lemma~\ref{lem:cauchy}. Ordering parameter blocks
as $d,p_1,\ldots,p_k$ and output groups in the corresponding
order makes the full Jacobian block lower triangular and
invertible. The inverse function theorem supplies local
isolation in this affine family, both for the selected-output
fiber and hence for the full-cube value fiber restricted to it.
\end{proof}

\paragraph{Scope for lower-bound counting.}
When also $\log H=O(n)$, these fibers eventually exceed
$H!C^{nH}$ for every fixed $C>0$, even though the representations
are counted as unordered. Thus such a universal bound on the
\emph{total} number of representations of one table is false.
When $2^{n-3k}\geq H$, the same cardinality counts isolated
regular solutions of the fixed selected-output constraints in
the specific $(k+1)H$-dimensional slice in
Theorem~\ref{thm:regular}. For example the growing-dimension
choice of $k$ satisfies this condition for
$H=\lfloor2^n/n\rfloor$ and all sufficiently large $n$, since
$2^{3k}\leq n^{3\log2/(2\log7)}$ and
$3\log2/(2\log7)<1$.
This still does not establish local isolation or full Jacobian
rank when all affine coefficients of every head vary independently.
Nor does it bound components after quotienting a larger continuous
fiber. It also does not lower bound the number of distinct
strict sign tables, since all the displayed representations compute
exactly the same real table.

\paragraph{The missing sign step.}
At any displayed representation, perturbing numerator and denominator
coefficients gives first-order directions in the span of
\[
 \frac{1}{B_h^\pi(x)},\quad
 \frac{x_i}{B_h^\pi(x)},\quad
 \frac{1}{(B_h^\pi(x))^2},\quad
 \frac{x_i}{(B_h^\pi(x))^2}.
\]
Subtracting one fixed representation in the fiber makes all the
centers identical to zero at a cost of another $H$ heads. A strict
sign vector attained by a tangent direction is attained by a small
actual perturbation, because the cube is finite and all denominator
slopes and values are strictly positive. However no statement that
these tangent spaces jointly meet all orthants is established.
The cardinality of the fiber by itself implies no such coverage.

\paragraph{Exact checks.}
The companion script enumerates all one-triple permutations for
$3\leq H\leq8$. The largest fibers have cardinalities
$2,3,6,12,28,76$, respectively. For two triples and $H\leq5$,
the complete joint signatures are injective. Thus separate
one-triple collisions cannot simply be combined independently:
mixed projection data can distinguish them. The script checks
the full-cube reciprocal identities with rational arithmetic
and checks all affine slopes directly. These finite results
test the construction and illustrate a constraint; the growing
fiber theorem uses the counting proof above.
\end{document}
